% ============================================================
% PREAMBLE BUKU AJAR KRIPTOGRAFI
% ============================================================

% ------------------------------------------------------------
% Encoding & Bahasa
% ------------------------------------------------------------
\usepackage[utf8]{inputenc}
\usepackage[T1]{fontenc}
\usepackage[indonesian]{babel}

% ------------------------------------------------------------
% Layout & Tipografi
% ------------------------------------------------------------
\usepackage[verbose=false]{geometry}
\geometry{
    a4paper,
    top=2.5cm,
    left=3cm,
    right=2cm,
    bottom=2cm
}

\setlength{\parindent}{0pt}
\setlength{\parskip}{0.5em}
\setlength{\emergencystretch}{2em}
\renewcommand{\baselinestretch}{1.5}
\raggedbottom
\usepackage{ragged2e}
\justifying

% ------------------------------------------------------------
% Paket Tabel & Struktur Dokumen
% ------------------------------------------------------------
\usepackage{booktabs}
\usepackage{longtable}
\usepackage{array}
\usepackage{enumitem}
\usepackage{subfiles}

% Nama job main document (untuk deteksi kompilasi per bab vs buku utuh)
\makeatletter
\def\MainDocJobname{buku_kriptografi}
\makeatother

% ------------------------------------------------------------
% Grafik, Diagram, & Hyperlink
% ------------------------------------------------------------
\usepackage{url}
\usepackage{graphicx}
\usepackage{xcolor}
\usepackage{tikz}
\usetikzlibrary{shapes.geometric, arrows.meta, positioning, calc}
\usepackage[hypertexnames=false]{hyperref}
\usepackage{bookmark}

% ------------------------------------------------------------
% Header & Footer
% ------------------------------------------------------------
\usepackage{fancyhdr}

% ------------------------------------------------------------
% Pengaturan Section (judul rata kiri, isi justified)
% ------------------------------------------------------------
\usepackage{titlesec}

\titleformat{\section}
  {\normalfont\Large\bfseries\raggedright}
  {\thesection}{1em}{}

\titleformat{\subsection}
  {\normalfont\large\bfseries\raggedright}
  {\thesubsection}{1em}{}

\titleformat{\subsubsection}
  {\normalfont\normalsize\bfseries\raggedright}
  {\thesubsubsection}{1em}{}

\titlespacing*{\section}{0pt}{3.5ex plus 1ex minus .2ex}{2.3ex plus .2ex}

% Pastikan isi paragraf tetap justified setelah judul
\usepackage{etoolbox}
\pretocmd{\section}{\justifying}{}{}
\pretocmd{\subsection}{\justifying}{}{}
\pretocmd{\subsubsection}{\justifying}{}{}

% ------------------------------------------------------------
% LISTING (Kode Program)
% ------------------------------------------------------------
\usepackage{listings}

\lstset{
    basicstyle=\ttfamily\small,
    breaklines=true,
    frame=single,
    numbers=left,
    numberstyle=\tiny,
    keywordstyle=\color{blue},
    commentstyle=\color{green!60!black},
    stringstyle=\color{red},
    showstringspaces=false,
    tabsize=2,
    float=false
}

\lstdefinestyle{matlab}{
    language=Matlab,
    morekeywords={function,end,if,else,for,while,return}
}

\lstdefinestyle{cstyle}{
    language=C,
    morekeywords={int,char,void,return,for,while,if,else}
}

\lstdefinestyle{cppstyle}{
    language=C++,
    morekeywords={int,char,void,return,for,while,if,else,string,cout,cin}
}

\lstdefinestyle{pascalstyle}{
    language=Pascal,
    morekeywords={program,var,begin,end,function,procedure,if,then,else,for,to,do,while}
}

% ------------------------------------------------------------
% KONTROL FLOAT (Gambar, Tabel, Listing)
% ------------------------------------------------------------
\usepackage{float}      % opsi [H]
\usepackage{placeins}   % \FloatBarrier

% Parameter float (lebih stabil untuk buku ajar)
\renewcommand{\topfraction}{0.9}
\renewcommand{\bottomfraction}{0.8}
\renewcommand{\textfraction}{0.05}
\renewcommand{\floatpagefraction}{0.85}
\setcounter{topnumber}{3}
\setcounter{bottomnumber}{2}
\setcounter{totalnumber}{4}

% Default posisi figure & table
\makeatletter
\def\fps@figure{!htbp}
\def\fps@table{!htbp}
\makeatother

% Cegah float melewati section
\let\oldsection\section
\renewcommand{\section}{\FloatBarrier\oldsection}

\let\oldsubsection\subsection
\renewcommand{\subsection}{\FloatBarrier\oldsubsection}

% Pastikan float selesai sebelum listing
\pretocmd{\lstlisting}{\FloatBarrier}{}{}

% ------------------------------------------------------------
% Metadata Dokumen
% ------------------------------------------------------------
\title{\textbf{Buku Ajar Kriptografi}\\
\large Untuk Program Studi Matematika S1}
\author{Departemen Matematika}
\date{}
